\section{Related Work}
\label{sec:related}

\paragraph{Symbolic and formal verification of mathematics.}
Computer-algebra systems such as SymPy \citep{meurer2017sympy} support exact
symbolic equivalence checking and are widely used as lightweight verifiers for
mathematical expressions. A complementary line of work aims at
\emph{formal} verification via proof assistants and autoformalization, translating
natural-language mathematics into machine-checkable statements
\citep{wu2022autoformalization}. These approaches are sound where they apply,
but they are limited to the fragment that can be parsed, formalized, or proven.
Our symbolic grader sits in the lightweight CAS regime: fast and exact on
covered steps, silent or inapplicable elsewhere.

\paragraph{LLM-as-judge and its failure modes.}
Using a strong LLM to grade another model's output is now a standard evaluation
tool \citep{zheng2023judging,liu2023geval}. Documented failure modes include
positional and other fairness biases \citep{wang2024faireval}, self-preference
for an evaluator's own generations \citep{panickssery2024selfpref}, and brittle
agreement with humans under adversarial or carefully curated pairs
\citep{zeng2024llmbar,thakur2025judging}. We study a related but distinct
failure in the math setting: judges that over-flag \emph{equivalence-preserving}
rewrites as errors, even when a symbolic checker correctly accepts them. By
running two judges from independent labs under one prompt and metric suite, we
separate model-independent patterns from model-specific operating points, and
reduce (though do not eliminate) the risk that a single judge's self-preference
dominates the comparison \citep{panickssery2024selfpref}.

\paragraph{Evaluating mathematical reasoning and locating errors.}
Math benchmarks such as GSM8K \citep{cobbe2021gsm8k} and MATH
\citep{hendrycks2021math} typically score final answers. Process-level
supervision and step-wise reward models push evaluation toward intermediate
reasoning \citep{uesato2022process,lightman2024verify,wang2024mathshepherd}.
Complementary \emph{step-quality} suites score chains for validity,
informativeness, or embedding-based consistency
\citep{golovneva2023roscoe,prasad2023receval,xia2025reasoneval}; ReasonEval in
particular evaluates error detection under perturbation, closest in spirit to
our injection suite \citep{xia2025reasoneval}. Meta-reasoning benchmarks such as
{MR-GSM8K} ask models to locate and explain mistakes in solutions
\citep{zeng2025mrgsm8k}. Closest to our human-set findings, \citet{tyen2024errors}
show that LLMs struggle to \emph{find} reasoning
errors even when they can \emph{correct} them given the error location: the
same bottleneck our judges hit on real, mostly out-of-fragment mistakes (F1) and
over-flagging of surface edits (F2). Rather than proposing a new grader or
benchmark, we evaluate two existing graders directly, scoring a symbolic
verifier and two LLM judges on the same multi-step derivations under one
protocol.

\paragraph{Abstention and selective prediction.}
Our framing, coverage vs.\ precision with abstention as a first-class
outcome, follows the tradition of reject options and selective
classification \citep{chow1970reject,elyaniv2010selective,geifman2017selective}
and its NLP instantiation as selective question answering
\citep{kamath2020selective}. We treat symbolic \textsc{none} and
confidence-thresholded judge abstention as two concrete reject policies for
math grading, and report coverage/precision trade-offs rather than forced
pass/fail alone.

\paragraph{Answer equivalence.}
Token-level match underestimates true answer quality when paraphrases and
semantically equivalent forms are treated as errors
\citep{bulian2022tomayto}. Our F2 finding is the process-level analogue:
equivalence-preserving algebraic rewrites are correct under a CAS checker but
often flagged by LLM judges.
